\documentclass[12pt,a4paper]{article}

%=========================================================
%                   ENCODAGE ET POLICES
%=========================================================
\usepackage[utf8]{inputenc}
\usepackage[T1]{fontenc}
\usepackage[french]{babel}
\usepackage{lmodern}

%=========================================================
%                     MATHÉMATIQUES
%=========================================================
\usepackage{amsmath,amssymb,amsfonts,mathrsfs}
\usepackage{tkz-tab}
\everymath{\displaystyle}

%=========================================================
%                     MISE EN PAGE
%=========================================================
\usepackage[left=1.5cm,right=1.5cm,top=1.5cm,bottom=1.8cm]{geometry}
\usepackage{multicol}
\usepackage{float}
\usepackage{array,multirow}
\usepackage{enumitem}
\usepackage{color, soul}

%=========================================================
%                GRAPHISMES ET COULEURS
%=========================================================
\usepackage{tikz}
\usetikzlibrary{arrows.meta,calc}
\usepackage{pgfplots}
\pgfplotsset{compat=1.15}

\usepackage[most]{tcolorbox}
\tcbuselibrary{skins,hooks}

%=========================================================
%                  LIENS ET ARRIÈRE-PLAN
%=========================================================
\usepackage[colorlinks=true,linkcolor=blue,urlcolor=blue,citecolor=blue]{hyperref}
\usepackage{eso-pic}

%=========================================================
%                   COULEURS ET STYLES
%=========================================================
\definecolor{myred}{RGB}{180, 40, 40}
\definecolor{myblue}{RGB}{20, 50, 120}
\definecolor{mygreen}{RGB}{30, 120, 40}

%=========================================================
%                       FILIGRANE
%=========================================================
\AddToShipoutPicture{
    \AtTextCenter{
        \makebox[0pt]{
            \rotatebox{45}{
                \textcolor[gray]{0.92}{
                    \fontsize{5cm}{5cm}\selectfont PGB
                }
            }
        }
    }
}

\begin{document}

% En-tête
\begin{center}
    \Large\textbf{\underline{\textcolor{myred}{Dénombrement}}}\\[0.2cm]
    \normalsize\textbf{Prof : M. BA} \hfill \textbf{Classe : Terminale L}\\[0.1cm]
    \textbf{Durée : 6 heures} \hfill \textbf{Année scolaire : 2025 -- 2026}
\end{center}
\hrule
\vspace{0.3cm}


\vfill
\hrule
\vspace{0.5cm}
\begin{center}
    \textit{« Compter n'est pas seulement dénombrer, c'est aussi organiser. »}
\end{center}
\vspace{0.5cm}
\hrule
\vfill



%=========================================================
%    I. Principes de base
%=========================================================
\section*{\color{myred}I. Principes fondamentaux du comptage}

\subsection*{\color{myred}1. Définition}
\begin{tcolorbox}[colback=red!5!white,colframe=myred,title=\textbf{Définition}]
Le \textbf{dénombrement} est l'ensemble des techniques mathématiques qui permettent de compter le nombre d'éléments d'un ensemble fini ou le nombre de possibilités d'une situation.
\end{tcolorbox}

\subsection*{\color{myred}2. Principe multiplicatif (Le produit)}

\begin{tcolorbox}[colback=blue!5!white,colframe=myblue,title=\textbf{Principe multiplicatif}]
Si une expérience comporte plusieurs étapes successives, le nombre total de possibilités est le \textbf{produit} du nombre de choix à chaque étape.
\end{tcolorbox}

\paragraph{Exemple :}
Un élève possède $3$ pantalons et $4$ chemises. Combien de tenues différentes peut-il composer ?
\begin{itemize}[label=$\bullet$]
    \item Étape 1 : Choisir un pantalon ($3$ choix).
    \item Étape 2 : Choisir une chemise ($4$ choix).
    \item Nombre total de tenues : $3 \times 4 = 12$ tenues.
\end{itemize}

\begin{tcolorbox}[colback=yellow!5!white,colframe=myred,title=\textbf{Exercice d'application 1}]
Un code de téléphone est composé de $4$ chiffres choisis parmi les chiffres $0, 1, 2, 3, 4, 5, 6, 7, 8, 9$.
\begin{enumerate}
    \item Combien de codes différents peut-on former si les répétitions sont autorisées ?
    \item Combien de codes peut-on former si tous les chiffres doivent être distincts ?
\end{enumerate}
\end{tcolorbox}


%=========================================================
%    II. Factorielle d'un entier
%=========================================================
\section*{\color{myred}II. Notion de factorielle}

\begin{tcolorbox}[colback=red!5!white,colframe=myred,title=\textbf{Définition}]
Pour tout entier naturel $n \ge 1$, on appelle \textbf{factorielle $n$}, notée $n!$, le produit de tous les entiers de $1$ jusqu'à $n$ :
\[
n! = n \times (n-1) \times \dots \times 2 \times 1
\]
Par convention, on pose : $\mathbf{0! = 1}$.
\end{tcolorbox}

\paragraph{Exemples :}
\begin{itemize}[label=$\bullet$]
    \item $1! = 1$
    \item $2! = 2 \times 1 = 2$
    \item $3! = 3 \times 2 \times 1 = 6$
    \item $4! = 4 \times 3 \times 2 \times 1 = 24$
    \item $5! = 5 \times 4 \times 3 \times 2 \times 1 = 120$
\end{itemize}

\end{document}